\section{Main Result}

\subsection{Theorem 2.1: Classical Closure Warrant}

The main theorem surface is the original incompressible Navier--Stokes
evolution.  The theorem below is aligned with the installed class-membership
route; it is not a substitute Euclidean, regularized, curvature, or
projected-flow claim.

\begin{theorem}[Classical closure warrant]
Let \(u_0\) be smooth divergence-free data for the three-dimensional
incompressible Navier--Stokes equation on the active theorem surface, and let
\((u,p)\) be the corresponding maximal classical solution on \([0,T_*)\).
Assume the installed same-fluid class-membership route:
\[
\begin{aligned}
RSCB.NKF&\Rightarrow NKF.Native\Rightarrow ACT.KX
\Rightarrow ACT.X\text{-Readout}\\
&\Rightarrow ACT.A\Rightarrow RCF.A\Rightarrow LCI.A
\Rightarrow FCI.5f\\
&\Rightarrow OFP.A\Rightarrow CFI.A\Rightarrow End_{NS}.
\end{aligned}
\]
Assume also the original-data endpoint admission package
\[
\begin{aligned}
&RawData.Production+EndpointFace.Type\\
&\quad+PackGauge.Dichotomy+RetainedAdmission.Persistence\\
&\Rightarrow H1+H2+H6+ORIGIN.Retain.
\end{aligned}
\]
Then no finite first class-membership endpoint occurs on the active
same-fluid route. Equivalently, every alleged finite terminal failure of the
retained class-membership packet is typed into
\[
Dead\vee packing\text{-}detached\vee tower\text{-}blown\vee Jump,
\]
and \(End_{NS}\) removes all four endpoint faces.
\end{theorem}

\begin{proof}
The retained receiver spine supplies \(RSCB.NKF\), hence \(NKF.Native\), and
therefore the \(ACT.KX\) differential budget on the retained same-fluid cover.
The installed readout and receiver steps give
\[
ACT.KX\Rightarrow ACT.X\text{-Readout}\Rightarrow ACT.A\Rightarrow RCF.A
\Rightarrow LCI.A.
\]
The source ledger is post-\(LCI.A\), so
\[
LCI.A+FCI.5f\Rightarrow CSP.A\Rightarrow OFP.A.
\]
Together with the class-entry side lemmas, \(OFP.A\) gives \(CFI.A\).

The endpoint package supplies
\[
\begin{aligned}
DTC\text{-}to\text{-}TowerBound
&+END.Exh+END.Cross\\
&+END.Pack+END.Field
\Rightarrow End_{NS}.
\end{aligned}
\]
The \(H1/H2/H6\) upgrade supplies the original-data endpoint admission: retained
native forcing persists unless an endpoint face appears, bounded pack gauge
persists unless the pack/tower face appears, and endpoint exhaustiveness is the
definition-level implication
\[
\begin{aligned}
\neg Part&\Rightarrow Dead,\\
\neg Pack&\Rightarrow packing\text{-}detached\vee tower\text{-}blown,\\
\neg Field&\Rightarrow Jump.
\end{aligned}
\]
Thus a first finite class-membership endpoint would have to be one of the four
listed faces, while \(End_{NS}\) excludes each face. This contradiction proves
the theorem.
\end{proof}

\paragraph{Claim boundary.}
This section records the warrant for the installed class-membership closure
route. It does not claim an export-ready Clay proof independently of the live
route surfaces, review surfaces, and release ledgers.
